\documentclass[10pt,a4paper]{amsart}
\usepackage[margin=19mm]{geometry}
\usepackage{amsmath,amssymb,mathtools}
\usepackage[scr=rsfs,cal=euler]{mathalfa}
\usepackage{xfrac}
\usepackage[unicode,hidelinks,bookmarks=false]{hyperref}
\newcommand{\ie}{i.e.\ }
\newcommand{\cf}{cf.\ }
\newcommand{\resp}{respectively\ }
\newcommand{\Subsection}{\subsection}
\providecommand{\nobreakdash}{\nobreak\hbox{-}}
\setlength{\emergencystretch}{2em}
\multlinegap0pt
\usepackage{amssymb}
\usepackage{bm}
\usepackage{float}
\usepackage{algorithm}\usepackage{algpseudocode}
\usepackage{xcolor}
\newtheorem{assumption}{Assumption}
\newtheorem{lemma}{Lemma}
\title{Strategic Index Reconstitution: Differential Games, Closed-Loop Equilibria and Mean-Field Dynamics}
\author{Lukas-Benedikt Fiechtner, Jose Blanchet}
\date{Source: arXiv:2609.15901v1 (CC BY 4.0)}
\begin{document}
\maketitle

\noindent\emph{Source and licence.} Strategic Index Reconstitution: Differential Games, Closed-Loop Equilibria and Mean-Field Dynamics. Version arXiv:2609.15901v1 (CC BY 4.0), distributed by arXiv under the Creative Commons Attribution 4.0 International licence, http://creativecommons.org/licenses/by/4.0/, as recorded in the arXiv licence metadata for this exact version and shown on the abstract page https://arxiv.org/abs/2609.15901v1. Full licence notice: \texttt{licenses/arxiv-2609.15901-CC-BY-4.0.txt}.\par\smallskip

\noindent\textit{Editorial scope.} Counted unit: the article's lemma lem:mf\_state\_transition\_representation (source0.tex lines 2823--2888). The article's complete original proof of it is delivered with it (source0.tex lines 4644--4807, one \texttt{proof} environment of 164 lines, which the article places away from the statement). No mathematics is written, repaired or re-derived: the statement and the proof are the article's own lines, in the article's own notation, and no retained line is substituted or re-flowed.  Retained uncounted as the source-local context the proof needs: lines 1205--1212; lines 2498--2504; lines 2542--2554; lines 2816--2818.  Nothing was omitted from any retained span, so the package carries no omission record.  No retained line contains a citation, so the wrapper carries no bibliography and no bibliography entry.  No retained line contains a figure environment, an image-inclusion command or drawing code, so no image is delivered and the package declares no asset dependency.  Editorial formatting only, in addition: an AMS \texttt{amsart} preamble that restates the article's own declarations and macros used by the retained lines, namely the environments \texttt{assumption}, \texttt{lemma} on separate counters, together with the standard packages those lines need.  The retained environments print in this document as Assumption 1, Lemma 1 in order of appearance, whereas the article prints the same items with the numbering of its own full document.  The complete native source of the article as distributed in the arXiv e-print for this exact version is bundled unchanged as \texttt{sources/210464/source0.tex}.
\begin{assumption}[Standing Assumptions]\label{ass:standing_model}
Under every \(\mathbb P_n^i\), \(S^0\) is a square-integrable càdlàg martingale
independent of \(\Xi\). Moreover, \(\Lambda\in\mathbb S_{++}^d\),
\(Q,Q_T\in\mathbb S_+^d\), \(\mathcal C(v)\geq0\) for every deterministic
\(v\in L^2([0,T],\mathbb R^d)\), \(\chi_u,\chi_g\geq0\), and
\(\varpi_n>0\). In each finite-player game, assume
\(\chi_u\varpi_n\leq1\).
\end{assumption}

\begin{equation}\label{eq:mf_riccati}
\dot P(t)
=
- Q+\frac12P(t)\Lambda^{-1}P(t),
\qquad
P(T)=Q_T.
\end{equation}

\begin{equation}\label{eq:mf_A_matrix}
\mathcal A(t)
=
\frac{1}{2+\varpi\chi_u}
\begin{pmatrix}
-\Lambda^{-1}P(t)&-\Lambda^{-1}&-\Lambda^{-1}\\
-\varpi\Gamma\Lambda^{-1}P(t)&
-(2+\varpi\chi_u)R-\varpi\Gamma\Lambda^{-1}&
-\varpi\Gamma\Lambda^{-1}\\
-\frac{\varpi\chi_u}{2}P(t)\Lambda^{-1}P(t)&
P(t)\Lambda^{-1}&P(t)\Lambda^{-1}
\end{pmatrix},
\end{equation}

\begin{equation}\label{eq:mf_matching_matrix}
\mathcal M:=\mathcal E_{\ell\ell}(T,0).
\end{equation}

\begin{lemma}
\label{lem:mf_state_transition_representation}
Under Assumption~\ref{ass:standing_model}, if \(\mathcal M\) is invertible,
the initial affine coefficient is
\begin{equation}\label{eq:mf_pre_matching}
\ell^0(0)
=
-\mathcal M^{-1}\sum_{k=1}^m\bar\pi^k\mathcal D_\ell^k.
\end{equation}
The pre-announcement solution is then
\begin{equation}\label{eq:mf_pre_forward_solution}
Y^0(t)
=
\mathcal E(t,0)
\begin{pmatrix}
0\\
0\\
\ell^0(0)
\end{pmatrix},
\qquad
t\in[0,T_{\mathrm{ann}}].
\end{equation}
In particular,
\[
\begin{aligned}
\bar X^0(T_{\mathrm{ann}})
=\mathcal E^{\mathrm{pre}}_{X\ell}\ell^0(0),
\qquad
I^0(T_{\mathrm{ann}})
=\mathcal E^{\mathrm{pre}}_{I\ell}\ell^0(0),
\qquad
\ell^0(T_{\mathrm{ann}})
=\mathcal E^{\mathrm{pre}}_{\ell\ell}\ell^0(0).
\end{aligned}
\]
If \(\mathcal E^{\mathrm{post}}_{\ell\ell}\) is invertible, then, using these
announcement values, the post-announcement initial affine coefficient for
scenario \(k\) is
\begin{equation}\label{eq:mf_post_transition_ell}
\ell^k(T_{\mathrm{ann}})
=
-(\mathcal E^{\mathrm{post}}_{\ell\ell})^{-1}
\left(\mathcal E^{\mathrm{post}}_{\ell X}\bar X^0(T_{\mathrm{ann}})+
\mathcal E^{\mathrm{post}}_{\ell I}I^0(T_{\mathrm{ann}})+\mathcal D_\ell^k\right),
\qquad k=1,\ldots,m.
\end{equation}
Finally, the post-announcement solution in scenario \(k\) is
\begin{equation}\label{eq:mf_post_forward_solution}
Y^k(t)
=
\mathcal E(t,T_{\mathrm{ann}})
\begin{pmatrix}
\bar X^0(T_{\mathrm{ann}})\\
I^0(T_{\mathrm{ann}})\\
\ell^k(T_{\mathrm{ann}})
\end{pmatrix}
+
\int_{T_{\mathrm{ann}}}^t
\mathcal E(t,s)b^k(s)\,ds,
\qquad
t\in[T_{\mathrm{ann}},T].
\end{equation}
If both \(\mathcal M\) and \(\mathcal E^{\mathrm{post}}_{\ell\ell}\) are
invertible, these formulas determine a unique solution of the mean-field
boundary value problem.
\end{lemma}

\begin{proof}
We first prove that the post-announcement block
\(\mathcal E^{\mathrm{post}}_{\ell\ell}\) is invertible. Recall from
Lemma~\ref{lem:mf_state_transition_representation} that
\(\mathcal E^{\mathrm{post}}=\mathcal E(T,T_{\mathrm{ann}})\), and that
\(\mathcal E^{\mathrm{post}}_{\ell\ell}\) denotes its \(\ell\ell\)-block. The matrix
\(\mathcal E(t,s)\) is the fundamental solution of the homogeneous ODE
\begin{equation}\label{eq:mf_homogeneous_transition_ode}
\dot{Y}(t)=\mathcal A(t)Y(t),
\qquad
Y(t)=(\bar X(t),I(t),\ell(t))^\top
\end{equation}
with \(\mathcal{A}(t)\) defined in \eqref{eq:mf_A_matrix}.
Let \(\mathfrak q\in\ker(\mathcal E^{\mathrm{post}}_{\ell\ell})\). Solve
\eqref{eq:mf_homogeneous_transition_ode} on
\([T_{\mathrm{ann}},T]\) with initial condition
\[
\bar X(T_{\mathrm{ann}})=0,
\qquad
I(T_{\mathrm{ann}})=0,
\qquad
\ell(T_{\mathrm{ann}})=\mathfrak q.
\]
By definition of \(\mathcal E^{\mathrm{post}}\) and the assumption
\(\mathfrak q\in\ker(\mathcal E^{\mathrm{post}}_{\ell\ell})\), the terminal affine coefficient is
\[
\begin{aligned}
\ell(T)
={}&
\mathcal E^{\mathrm{post}}_{\ell X}\bar X(T_{\mathrm{ann}})
+
\mathcal E^{\mathrm{post}}_{\ell I}I(T_{\mathrm{ann}})
+
\mathcal E^{\mathrm{post}}_{\ell\ell}\ell(T_{\mathrm{ann}})
=\mathcal E^{\mathrm{post}}_{\ell\ell}\mathfrak q
=0.
\end{aligned}
\]
The first row of \eqref{eq:mf_homogeneous_transition_ode} is
\begin{equation}\label{eq:mf_homogeneous_mean_inventory_ode}
\dot{\bar X}(t)
=
-\frac{1}{2+\varpi\chi_u}\Lambda^{-1}
\left(P(t)\bar X(t)+I(t)+\ell(t)\right).
\end{equation}
Denote the right-hand side of
\eqref{eq:mf_homogeneous_mean_inventory_ode} by \(\bar u(t)\) and define the
costate \(p(t):=P(t)\bar X(t)+\ell(t)\). Then
\eqref{eq:mf_homogeneous_mean_inventory_ode} can be reformulated as
\begin{equation}\label{eq:mf_homogeneous_first_row_identity}
p(t)+I(t)+(2+\varpi\chi_u)\Lambda\bar u(t)=0.
\end{equation}
The costate process \(p(t)\) satisfies
\[
\dot p(t)=- Q\bar X(t),
\qquad
p(T)=Q_T\bar X(T).
\]
Indeed, the terminal condition follows from \(P(T)=Q_T\) and \(\ell(T)=0\). The ODE
follows by differentiating \(p(t)=P(t)\bar X(t)+\ell(t)\) and using
\eqref{eq:mf_riccati} together with the first and third rows of
\eqref{eq:mf_homogeneous_transition_ode}.
The impact row of \eqref{eq:mf_homogeneous_transition_ode} gives
\begin{equation}\label{eq:mf_homogeneous_impact_row}
\dot I(t)=-RI(t)+\varpi\Gamma\bar u(t).
\end{equation}
Multiplying \eqref{eq:mf_homogeneous_first_row_identity} by
\(\bar u(t)^\top\) and
integrating over \([T_{\mathrm{ann}},T]\) gives
\begin{equation}\label{eq:mf_homogeneous_energy_identity}
0
=
\int_{T_{\mathrm{ann}}}^T\bar u(t)^\top p(t)\,dt
+
\int_{T_{\mathrm{ann}}}^T\bar u(t)^\top I(t)\,dt
+
(2+\varpi\chi_u)\int_{T_{\mathrm{ann}}}^T
\bar u(t)^\top\Lambda\bar u(t)\,dt.
\end{equation}
Integration by parts gives
\[
\begin{aligned}
\int_{T_{\mathrm{ann}}}^T\bar u(t)^\top p(t)\,dt
={}&\bar X(T)^\top p(T)
-\bar X(T_{\mathrm{ann}})^\top p(T_{\mathrm{ann}})-\int_{T_{\mathrm{ann}}}^T\bar X(t)^\top\dot p(t)\,dt\\
={}&\bar X(T)^\top Q_T\bar X(T)
+\int_{T_{\mathrm{ann}}}^T\bar X(t)^\top Q\bar X(t)\,dt.
\end{aligned}
\]
All three summands on the right-hand side of
\eqref{eq:mf_homogeneous_energy_identity} are nonnegative. The first is
nonnegative by the integration-by-parts identity since
\(Q,Q_T\succeq0\). The second is nonnegative because
\eqref{eq:mf_homogeneous_impact_row}, together with
\(I(T_{\mathrm{ann}})=0\), identifies it with \(\varpi\) times the transient-impact quadratic
form of \(\bar u\) on \([T_{\mathrm{ann}},T]\), which is nonnegative by
\(\varpi>0\) and the
kernel-positivity requirement in Assumption~\ref{ass:standing_model}. The third
is nonnegative by \(2+\varpi\chi_u>0\) and \(\Lambda\succ0\). Hence
\eqref{eq:mf_homogeneous_energy_identity} implies \(\bar u\equiv0\). Since
\(\bar X(T_{\mathrm{ann}})=0\) and
\(\dot{\bar X}(t)=\bar u(t)\), this implies \(\bar X\equiv0\). The impact row
\eqref{eq:mf_homogeneous_impact_row} with \(I(T_{\mathrm{ann}})=0\) then
gives \(I\equiv0\). The identity \(p+I+(2+\varpi\chi_u)\Lambda\bar u=0\) from
\eqref{eq:mf_homogeneous_first_row_identity} gives
\(p\equiv0\). Evaluating at \(T_{\mathrm{ann}}\) therefore yields
\[
0=p(T_{\mathrm{ann}})
=P(T_{\mathrm{ann}})\bar X(T_{\mathrm{ann}})+\ell(T_{\mathrm{ann}})
=\ell(T_{\mathrm{ann}})=\mathfrak q
\]
and so \(\mathfrak q=0\). Hence
\(\ker(\mathcal E^{\mathrm{post}}_{\ell\ell})=\{0\}\) and
\(\mathcal E^{\mathrm{post}}_{\ell\ell}\) is invertible.

It remains to prove that the matching matrix \(\mathcal M\)
from \eqref{eq:mf_matching_matrix} is invertible.
Let \(\mathfrak q\in\ker(\mathcal M)\). Run
\eqref{eq:mf_homogeneous_transition_ode} on \([0,T_{\mathrm{ann}}]\) from
\[
\bar X(0)=0,
\qquad
I(0)=0,
\qquad
\ell(0)=\mathfrak q.
\]
By definition of \(\mathcal E^{\mathrm{pre}}\), the state at \(T_{\mathrm{ann}}\) is
\[
\bar X(T_{\mathrm{ann}})=\mathcal E^{\mathrm{pre}}_{X\ell}\mathfrak q,
\qquad
I(T_{\mathrm{ann}})=\mathcal E^{\mathrm{pre}}_{I\ell}\mathfrak q,
\qquad
\ell(T_{\mathrm{ann}})=\mathcal E^{\mathrm{pre}}_{\ell\ell}\mathfrak q.
\]
Using
\(\mathcal E(T,0)=\mathcal E(T,T_{\mathrm{ann}})
\mathcal E(T_{\mathrm{ann}},0)\) and
\eqref{eq:mf_matching_matrix}, the assumption
\(\mathfrak q\in\ker(\mathcal M)\) gives
\[
\mathcal E^{\mathrm{post}}_{\ell X}\bar X(T_{\mathrm{ann}})
+
\mathcal E^{\mathrm{post}}_{\ell I}I(T_{\mathrm{ann}})
+
\mathcal E^{\mathrm{post}}_{\ell\ell}\ell(T_{\mathrm{ann}})
=\mathcal M\mathfrak q
=0.
\]
By definition of \(\mathcal E^{\mathrm{post}}\), this implies that the solution to
\eqref{eq:mf_homogeneous_transition_ode} on
\([T_{\mathrm{ann}},T]\), starting from
\((\bar X(T_{\mathrm{ann}}),I(T_{\mathrm{ann}}),\ell(T_{\mathrm{ann}}))\), has terminal
affine coefficient \(\ell(T)=0\). Equivalently, running the homogeneous ODE over
the combined interval \([0,T]\) with \(\bar X(0)=0\), \(I(0)=0\), and
\(\ell(0)=\mathfrak q\)
produces a solution satisfying \(\ell(T)=0\). The same energy argument used for
\(\mathcal E^{\mathrm{post}}_{\ell\ell}\), now on \([0,T]\), gives
\(\bar u\equiv0\), hence \(\bar X\equiv0\), \(I\equiv0\), and
\(p\equiv0\).
Since \(p\equiv0\), \(\bar X(0)=0\), and \(\ell(0)=\mathfrak q\), evaluating
\(p=P\bar X+\ell\) at \(t=0\) gives
\(0=p(0)=\mathfrak q\). Thus
\(\ker(\mathcal M)=\{0\}\), and \(\mathcal M\) is invertible.
\end{proof}

\end{document}
